% Copyright (c) 2026 Geovana Neves. Licensed under CC BY 4.0 (https://creativecommons.org/licenses/by/4.0/). See LICENSE-AND-AUTHORSHIP.md.
%
% 04-solver-setup/main.tex
%
% Guide 4 of the fts-workspace guides: Solver setup. Built by ../build-all.ps1
% (or build-all.sh) with this folder and ../shared/ as include directories.
%
% No em dashes and no en dashes. No underscores in frame or section titles
% outside \code: beamer writes titles to .nav and .toc.

\documentclass[aspectratio=169,11pt]{beamer}

\input{preamble-slides.tex}
\input{hooks-slides.tex}
\ftsinput{boxes}
\ftsinput{hooks-contract}

\usepackage[nohyphen]{underscore}

\input{info.tex}

\newcommand{\setupsources}[1]{%
  \par\vfill
  {\tiny\color{SoftGray} Sources under \code{src/pyflightstream/}: #1.\par}}

% Keep the slide edition's section-frame shape and add its source line.
\AtBeginSection[]{%
  \begin{frame}[plain,noframenumbering]
    \vfill
    {\color{SoftGray}\small Part \thesection\par}
    \vspace{2mm}
    {\color{hanblue}\Huge\bfseries\insertsection\par}
    \vfill
    \setupsources{\code{cases/\_\_init\_\_.py}, \code{SolverSettings};
      \code{workspace/matrix.py}; \code{cases/workflows.py}; \code{commands/}}
  \end{frame}}

\begin{document}

\guidetitleframe{4}{Solver setup}{Solver setup}
  {The setup file, its keys, and the solver commands behind them}
  {Source: the package page ``The workspace and the workflow'' and the setup code [1].}

\begin{frame}{Contents}
\small
\begin{multicols}{2}
\tableofcontents
\end{multicols}
\setupsources{\code{workspace/matrix.py}; \code{cases/workflows.py};
  \code{script/helpers.py}; \code{commands/advanced\_settings.yaml}}
\end{frame}

% Copyright (c) 2026 Geovana Neves. Licensed under CC BY 4.0 (https://creativecommons.org/licenses/by/4.0/). See LICENSE-AND-AUTHORSHIP.md.
%
\section{The setup artifact}

\begin{frame}{Where the solver setup lives}
\begin{itemize}\small
  \item The matrix row's \code{SET} cell names the setup artifact:
        \code{s001} selects \code{inputs/setups/s001.toml}.
        The general filename is \code{sNNN.toml}.
  \item Top-level TOML keys become \code{SolverSettings} fields, directly
        or through the preset alias table. A key that reaches nothing and
        is not declared recorded-only is refused at plan.
  \item An unstated optional setting normally emits nothing, leaving the
        solver's own default. The package also has explicit defaults and
        derived values, listed on the next frame.
  \item The run record keeps the \code{solver\_setup} snapshot: effective
        flag values and provenance, including defaults and unknown values.
        Unknown does not mean disabled.
\end{itemize}
\setupsources{\code{workspace/inputs.py}, \code{SetupArtifact};
  \code{workspace/matrix.py}, \code{\_solver\_from\_setup};
  \code{script/solver\_setup.py}, \code{SolverSetup}; \code{script/helpers.py}}
\end{frame}

\begin{frame}{Package defaults and solver defaults}
\begin{itemize}\small
  \item \code{iterations = 500} and \code{convergence = 1e-5} are model
        defaults and reach solver commands even when omitted from the file.
  \item Every other \code{SolverSettings} field defaults to Python
        \code{None}. In TOML, omit the key; there is no \code{None} value.
  \item An omitted \code{minimum\_cp} still emits
        \code{SOLVER\_MINIMUM\_CP -100} on a build carrying that command.
        The solver's documented default is \code{-20}.
  \item The workflow supplies \code{INCOMPRESSIBLE} for an unstated model,
        the freestream velocity for an unstated reference velocity, and
        \code{DISABLE} for loading saved solver initialization.
  \item The snapshot distinguishes explicit, default and unknown values.
        Its flag inventory is \code{FLAG\_SPECS}; that inventory is larger
        than the set of keys accepted in a setup file.
\end{itemize}
\setupsources{\code{cases/\_\_init\_\_.py}, \code{SolverSettings};
  \code{cases/workflows.py}; \code{script/helpers.py}, \code{solver\_settings};
  \code{script/solver\_setup.py}, \code{FLAG\_SPECS}}
\end{frame}

\begin{frame}{Recorded-only keys}
\relax
{\ftstablesize
\begin{tabularx}{\linewidth}{@{}>{\RaggedRight\arraybackslash}p{0.43\linewidth}>{\RaggedRight\arraybackslash}X@{}}
\toprule
\textbf{Key} & \textbf{Why it emits nothing} \\
\midrule
\code{solver}, \code{motion} & Run type and motion come from the row and workflow. \\
\code{symmetry\_type} & Geometry symmetry belongs to the row's \code{SYMMETRY}. \\
\code{unsteady\_delta\_theta\_deg} & Superseded by the row's \code{DELTA\_THETA}. \\
\code{unsteady\_N\_revolutions} & Superseded by the row's \code{REVOLUTIONS}. \\
\code{set\_base\_region\_trailing\_edges} & Boundary-selecting separation model: a recipe's job. \\
\code{slipstream\_wake\_stabilization}, \code{wake\_layers} & No emitter in this package. \\
\bottomrule
\end{tabularx}}
\begin{itemize}\small
  \item Each warns by name and reason. A file may declare additional names
        with \code{recorded\_only = ["my\_setting"]}; these still warn.
  \item Do not use recorded-only to make a misspelled setting appear to work.
        \code{mesh\_order\_list} is refused outright: boundary order belongs
        to the geometry inventory.
\end{itemize}
\setupsources{\code{workspace/matrix.py}, \code{\_PRESET\_RECORDED\_ONLY},
  \code{\_PRESET\_RECORDED\_ONLY\_KEY}, \code{\_solver\_from\_setup}}
\end{frame}

% Copyright (c) 2026 Geovana Neves. Licensed under CC BY 4.0 (https://creativecommons.org/licenses/by/4.0/). See LICENSE-AND-AUTHORSHIP.md.
%
\section{Run control}

\begin{frame}{Iteration and convergence controls}
\relax
{\ftstablesize
\begin{tabularx}{\linewidth}{@{}>{\RaggedRight\arraybackslash}p{0.25\linewidth}>{\RaggedRight\arraybackslash}p{0.38\linewidth}>{\RaggedRight\arraybackslash}X@{}}
\toprule
\textbf{Setup key} & \textbf{Solver command} & \textbf{Type and package default} \\
\midrule
\code{iterations} (alias \code{NITER}) & \code{SOLVER\_SET\_ITERATIONS} & Integer $\geq 1$; \code{500}. \\
\code{convergence} & \code{SOLVER\_SET\_CONVERGENCE} & Float $>0$; \code{1e-5}. \\
\code{forced\_iterations} & \code{SOLVER\_SET\_FORCED\_ITERATIONS} & Toggle; \code{None}. \\
\code{convergence\_iterations} & \code{SET\_SOLVER\_CONVERGENCE\_ITERATIONS} & Integer $\geq 1$; \code{None}. \\
\bottomrule
\end{tabularx}}
\begin{itemize}\small
  \item \code{iterations} is the solver iteration limit. Convergence is
        the residual threshold on surface velocity and pressure.
  \item Forced iterations run the full count regardless of convergence.
        Convergence iterations is how long the residual must remain below
        the threshold before convergence is declared.
  \item A toggle accepts \code{true}/\code{false} or the solver words
        \code{"ENABLE"}/\code{"DISABLE"}. Numbers must be finite.
\end{itemize}
\setupsources{\code{cases/\_\_init\_\_.py}, \code{SolverSettings};
  \code{workspace/matrix.py}, \code{\_PRESET\_ALIASES};
  \code{commands/runtime\_settings.yaml}; \code{commands/advanced\_settings.yaml}; \code{script/helpers.py}}
\end{frame}

\begin{frame}{Threads and wall-clock limits}
\begin{itemize}\small
  \item \code{max\_threads}: integer $\geq 1$, default \code{None}; alias
        \code{max\_parallel\_threads}. Parallel core count, emitted as
        \code{SET\_MAX\_PARALLEL\_THREADS}. A stated row \code{NCPUS}
        overrides the setup value; an unstated row inherits it.
  \item \code{timeout\_s}: finite float $>0$, default \code{None}.
        Wall-clock limit for one point's solver process, enforced by the
        executor, not by a FlightStream command.
  \item \code{walltime\_margin\_s}: finite float $>0$, default \code{None}.
        Seconds of the row's \code{WALLTIME} reserved for exports.
        When unstated, the wall-clock program uses 20 minutes.
  \item The margin belongs to the setup. It emits no solver command;
        the wall-clock program written beside the script reads it.
\end{itemize}
\setupsources{\code{cases/\_\_init\_\_.py}, \code{SolverSettings};
  \code{cases/workflows.py}, \code{row\_ncpus}; \code{workspace/matrix.py};
  \code{commands/runtime\_settings.yaml}}
\end{frame}

\begin{frame}{NITER: measured on build 26.124}
\begin{itemize}\small
  \item One short rotor case, 36 physical time steps: with \code{NITER = 5},
        every step used 5 inner iterations, 180 total; \code{COMPLETED\_MAX\_ITER}.
  \item With \code{NITER = 300} and \code{convergence\_iterations = 3},
        the case converged in 6 to 90 inner iterations per step,
        median 32.5, 1100 total.
  \item On this measured build and geometry, \code{NITER} caps each time step.
        Convergence iterations counts consecutive iterations satisfying the
        criterion; it is not an iteration cap.
\end{itemize}
\begin{block}{Budget and scope}
Worst-case inner work is the number of time steps times \code{NITER}.
This is one build and one geometry, not verification of every solver build.
\end{block}
{\tiny Release evidence: measurement record, check \code{niter\_limit} [2].\par}
\end{frame}

% Copyright (c) 2026 Geovana Neves. Licensed under CC BY 4.0 (https://creativecommons.org/licenses/by/4.0/). See LICENSE-AND-AUTHORSHIP.md.
%
\section{Physics and model}

\begin{frame}{Boundary layer and viscous coupling}
\begin{itemize}\small
  \item \code{boundary\_layer}: string, default \code{None}; alias
        \code{boundary\_layer\_type}. Emits \code{SET\_BOUNDARY\_LAYER\_TYPE}
        with \code{LAMINAR}, \code{TRANSITIONAL} or \code{TURBULENT}.
        The database gives \code{TRANSITIONAL} as the solver default.
  \item Its notes state the default model's chord Reynolds range as
        500000 to 1500000, with an explicit laminar or turbulent choice
        outside that range [3, p.203].
  \item \code{viscous\_coupling}: toggle, default \code{None}; emits
        \code{SET\_SOLVER\_VISCOUS\_COUPLING}. It couples the boundary
        layer model to the potential solution. Stall capture also needs
        a flow-separation model assignment.
  \item \code{reynolds\_averaged\_drag}: toggle, default \code{None}; alias
        \code{reynolds\_averaged\_drag\_forces}. Emits
        \code{REYNOLDS\_AVERAGED\_DRAG\_FORCES}, the flat-plate
        boundary-layer calculation toggle. No solver default is stated.
\end{itemize}
\setupsources{\code{cases/\_\_init\_\_.py}; \code{workspace/matrix.py};
  \code{cases/workflows.py}; \code{script/helpers.py};
  \code{commands/solver\_settings.yaml}; \code{commands/advanced\_settings.yaml}}
\end{frame}

\begin{frame}{Model and initialization}
\begin{itemize}\small
  \item \code{solver\_model}: string, default \code{None}; alias
        \code{set\_solver\_model}. The workflow passes it to
        \code{INITIALIZE\_SOLVER}, using \code{INCOMPRESSIBLE} if omitted.
  \item Accepted model arguments: \code{INCOMPRESSIBLE},
        \code{SUBSONIC\_PRANDTL\_GLAUERT}, \code{TRANSONIC\_FIELD\_PANEL},
        \code{TANGENT\_CONE}, \code{MODIFIED\_NEWTONIAN}.
  \item \code{wall\_collision\_avoidance}: toggle, default \code{None};
        alias \code{proximity\_avoidance}. Another argument of
        \code{INITIALIZE\_SOLVER}; the database says it applies to
        the first three models above.
  \item \code{load\_solver\_initialization}: toggle, default \code{None}.
        Reaches \code{OPEN}'s \code{LOAD\_SOLVER\_INITIALIZATION} argument.
        The workflow uses \code{DISABLE} if omitted, so a saved solver state
        is not loaded implicitly.
\end{itemize}
\setupsources{\code{cases/\_\_init\_\_.py}; \code{workspace/matrix.py};
  \code{cases/workflows.py}; \code{script/helpers.py}, \code{initialize\_solver};
  \code{commands/solver\_initialization.yaml}; \code{commands/file\_io.yaml}}
\end{frame}

\begin{frame}{Stabilization and pressure floor}
\begin{itemize}\small
  \item \code{solver\_stabilization}: finite float $\geq 0$, default
        \code{None}; emits \code{SOLVER\_STABILIZATION <value>}.
        The database describes a dimensionless strength from 0 to 1;
        the model does not impose that upper bound. No solver default is stated.
  \item Alternative preset form: \code{stabilization = "ENABLE"} and
        \code{stabilization\_strength = 0.25}. A strength without a gate
        is enabled. Enabling without a strength is refused.
  \item \code{stabilization = "DISABLE"} produces no stabilization line.
        Do not combine the gated pair with \code{solver\_stabilization}.
        Zero is an emitted numeric value, even though the manual describes
        zero strength as equivalent to disabling.
  \item \code{minimum\_cp}: finite float, default \code{None}; alias
        \code{solver\_minimum\_cp}. It sets the pressure-coefficient floor
        with \code{SOLVER\_MINIMUM\_CP}. Omission uses the library's
        \code{-100}, overriding the solver default \code{-20}.
\end{itemize}
\setupsources{\code{cases/\_\_init\_\_.py}; \code{workspace/matrix.py},
  \code{\_resolve\_stabilization}; \code{script/helpers.py};
  \code{commands/solver\_settings.yaml}; \code{commands/advanced\_settings.yaml}}
\end{frame}

\begin{frame}{Boundary families}
\begin{itemize}\small
  \item \code{vorticity\_drag\_families}: list of strings, default
        \code{None}; aliases \code{vorticity\_drag\_boundaries} and
        \code{set\_vorticity\_drag\_boundaries}. Reaches
        \code{SET\_VORTICITY\_DRAG\_BOUNDARIES} in analysis.
  \item \code{axial\_separation\_families}: list of strings, default
        \code{None}; no alias. Reaches
        \code{SET\_AXIAL\_SEPARATION\_BOUNDARIES} in init.
        Documented through 26.100; a later build is refused at plan.
  \item Write family names, not indices. Both lists resolve against the
        opened geometry: absent families are left out; a selection resolving
        to nothing is refused. An omitted key emits no selection or deletion.
  \item Without a vorticity selection, the solver uses surface-pressure
        integration for induced drag. The helper notes advise against a
        blanket vorticity selection on mixed wing and bluff-body geometry.
\end{itemize}
\setupsources{\code{cases/\_\_init\_\_.py}; \code{workspace/matrix.py};
  \code{cases/workflows.py}, \code{\_vorticity\_indices}, \code{\_axial\_separation\_indices};
  \code{script/helpers.py}; \code{commands/solver\_settings.yaml}}
\end{frame}

\begin{frame}{Reference velocity, precision and symmetry loads}
\begin{itemize}\small
  \item \code{reference\_velocity\_m\_per\_s}: finite float $>0$, default
        \code{None}; alias \code{reference\_velocity\_mps}. Emits
        \code{SOLVER\_SET\_REF\_VELOCITY}, in m/s, for coefficient
        normalization. Omission makes the workflow state freestream velocity.
  \item \code{significant\_digits}: integer $\geq 1$, default \code{None};
        emits \code{SET\_SIGNIFICANT\_DIGITS} in setup. The command database
        describes output precision from 1 to 8; the model has no upper bound.
        The model's comment records the solver default as 4 decimals.
  \item \code{symmetry\_loads}: toggle, default \code{None}; emits
        \code{SET\_ANALYSIS\_SYMMETRY\_LOADS} before solving.
        It moved to the row's \code{SYMMETRY\_LOADS} column, with the
        preset still accepted as a fallback.
  \item A stated row value wins; disagreement with the setup warns.
        True includes the whole wheel, false the meshed sector. If neither
        states it, no line is emitted.
\end{itemize}
\setupsources{\code{cases/\_\_init\_\_.py}; \code{workspace/matrix.py};
  \code{cases/workflows.py}, \code{\_row\_symmetry\_loads};
  \code{commands/runtime\_settings.yaml}; \code{commands/simulation\_controls.yaml};
  \code{commands/solver\_analysis.yaml}}
\end{frame}

% Copyright (c) 2026 Geovana Neves. Licensed under CC BY 4.0 (https://creativecommons.org/licenses/by/4.0/). See LICENSE-AND-AUTHORSHIP.md.
%
\section{Advanced settings with setup keys}

\begin{frame}{Convergence and pressure floor}
\small All commands here are inline, in phase \code{init}. Package default:
\code{None} for every key below. Floats must be finite.
\par\vspace{3mm}
{\ftstablesize\setlength{\tabcolsep}{3pt}
\begin{tabularx}{\linewidth}{@{}>{\RaggedRight\arraybackslash}p{0.18\linewidth}>{\RaggedRight\arraybackslash}p{0.18\linewidth}>{\RaggedRight\arraybackslash}p{0.26\linewidth}>{\RaggedRight\arraybackslash}p{0.07\linewidth}>{\RaggedRight\arraybackslash}p{0.10\linewidth}>{\RaggedRight\arraybackslash}X@{}}
\toprule
\textbf{Key} & \textbf{Alias} & \textbf{Solver command} & \textbf{Type} & \textbf{Solver default} & \textbf{26.123 / 26.124} \\
\midrule
\code{convergence\_\allowbreak iterations} & None & \code{SET\_\allowbreak SOLVER\_\allowbreak CONVERGENCE\_\allowbreak ITERATIONS} & int $\geq 1$ & Not stated & V / V \\
\code{minimum\_\allowbreak cp} & \code{solver\_\allowbreak minimum\_\allowbreak cp} & \code{SOLVER\_\allowbreak MINIMUM\_\allowbreak CP} & float & -20 & V / V \\
\bottomrule
\end{tabularx}}
\par\vspace{3mm}
\small \textbf{V}: verified on that build. \textbf{D}: documented, not verified on that build.
\par\vspace{2mm}
\small The first is the count below the residual threshold required for convergence. The second is the pressure-coefficient floor. An omitted minimum Cp emits the library default \code{-100}; the solver default is \code{-20}.
\par\vspace{2mm}{\scriptsize Manual references in the database: [3, p.344]; [3, p.345].}

\setupsources{\code{cases/\_\_init\_\_.py}, \code{SolverSettings}; \code{workspace/matrix.py}, \code{\_PRESET\_ALIASES}; \code{cases/workflows.py}; \code{script/helpers.py}; \code{commands/advanced\_settings.yaml}}
\end{frame}

\begin{frame}{Drag and mesh-induced wake velocity}
\small All commands here are inline, in phase \code{init}. Package default:
\code{None} for every key below. Floats must be finite.
\par\vspace{3mm}
{\ftstablesize\setlength{\tabcolsep}{3pt}
\begin{tabularx}{\linewidth}{@{}>{\RaggedRight\arraybackslash}p{0.18\linewidth}>{\RaggedRight\arraybackslash}p{0.18\linewidth}>{\RaggedRight\arraybackslash}p{0.26\linewidth}>{\RaggedRight\arraybackslash}p{0.07\linewidth}>{\RaggedRight\arraybackslash}p{0.10\linewidth}>{\RaggedRight\arraybackslash}X@{}}
\toprule
\textbf{Key} & \textbf{Alias} & \textbf{Solver command} & \textbf{Type} & \textbf{Solver default} & \textbf{26.123 / 26.124} \\
\midrule
\code{reynolds\_\allowbreak averaged\_\allowbreak drag} & \code{reynolds\_\allowbreak averaged\_\allowbreak drag\_\allowbreak forces} & \code{REYNOLDS\_\allowbreak AVERAGED\_\allowbreak DRAG\_\allowbreak FORCES} & toggle & Not stated & V / D \\
\code{mesh\_\allowbreak induced\_\allowbreak wake\_\allowbreak velocity} & \code{induced\_\allowbreak wake\_\allowbreak velocity} & \code{SOLVER\_\allowbreak SET\_\allowbreak MESH\_\allowbreak INDUCED\_\allowbreak WAKE\_\allowbreak VELOCITY} & toggle & Not stated & V / D \\
\bottomrule
\end{tabularx}}
\par\vspace{3mm}
\small \textbf{V}: verified on that build. \textbf{D}: documented, not verified on that build.
\par\vspace{2mm}
\small \code{reynolds\_averaged\_drag} switches the flat-plate boundary-layer calculations. For these toggles, the database gives no solver default; omission does not assert \code{DISABLE}.
\par\vspace{2mm}{\scriptsize Manual references in the database: [3, p.344].}

\setupsources{\code{cases/\_\_init\_\_.py}, \code{SolverSettings}; \code{workspace/matrix.py}, \code{\_PRESET\_ALIASES}; \code{cases/workflows.py}; \code{script/helpers.py}; \code{commands/advanced\_settings.yaml}}
\end{frame}

\begin{frame}{Farfield and unsteady pressure}
\small All commands here are inline, in phase \code{init}. Package default:
\code{None} for every key below. Floats must be finite.
\par\vspace{3mm}
{\ftstablesize\setlength{\tabcolsep}{3pt}
\begin{tabularx}{\linewidth}{@{}>{\RaggedRight\arraybackslash}p{0.18\linewidth}>{\RaggedRight\arraybackslash}p{0.18\linewidth}>{\RaggedRight\arraybackslash}p{0.26\linewidth}>{\RaggedRight\arraybackslash}p{0.07\linewidth}>{\RaggedRight\arraybackslash}p{0.10\linewidth}>{\RaggedRight\arraybackslash}X@{}}
\toprule
\textbf{Key} & \textbf{Alias} & \textbf{Solver command} & \textbf{Type} & \textbf{Solver default} & \textbf{26.123 / 26.124} \\
\midrule
\code{farfield\_\allowbreak layers} & None & \code{SOLVER\_\allowbreak SET\_\allowbreak FARFIELD\_\allowbreak LAYERS} & int 1 to 5 & 3 & V / D \\
\code{unsteady\_\allowbreak pressure\_\allowbreak and\_\allowbreak kutta} & \code{unsteady\_\allowbreak pressure\_\allowbreak kutta} & \code{SOLVER\_\allowbreak UNSTEADY\_\allowbreak PRESSURE\_\allowbreak AND\_\allowbreak KUTTA} & toggle & Not stated & V / D \\
\bottomrule
\end{tabularx}}
\par\vspace{3mm}
\small \textbf{V}: verified on that build. \textbf{D}: documented, not verified on that build.
\par\vspace{2mm}
\small The documented farfield layer range is \textbf{1 to 5}, solver default \textbf{3}; a value above 5 is refused, and every standard setup states \textbf{5}. The pressure toggle controls the unsteady Bernoulli and Kutta terms.
\par\vspace{2mm}{\scriptsize Manual references in the database: [3, p.344].}

\setupsources{\code{cases/\_\_init\_\_.py}, \code{SolverSettings}; \code{workspace/matrix.py}, \code{\_PRESET\_ALIASES}; \code{cases/workflows.py}; \code{script/helpers.py}; \code{commands/advanced\_settings.yaml}}
\end{frame}

\begin{frame}{Wake termination units}
\small All commands here are inline, in phase \code{init}. Package default:
\code{None} for every key below. Floats must be finite.
\par\vspace{3mm}
{\ftstablesize\setlength{\tabcolsep}{3pt}
\begin{tabularx}{\linewidth}{@{}>{\RaggedRight\arraybackslash}p{0.18\linewidth}>{\RaggedRight\arraybackslash}p{0.18\linewidth}>{\RaggedRight\arraybackslash}p{0.26\linewidth}>{\RaggedRight\arraybackslash}p{0.07\linewidth}>{\RaggedRight\arraybackslash}p{0.10\linewidth}>{\RaggedRight\arraybackslash}X@{}}
\toprule
\textbf{Key} & \textbf{Alias} & \textbf{Solver command} & \textbf{Type} & \textbf{Solver default} & \textbf{26.123 / 26.124} \\
\midrule
\code{wake\_\allowbreak termination\_\allowbreak steps} & None & \code{SET\_\allowbreak WAKE\_\allowbreak TERMINATION\_\allowbreak TIME\_\allowbreak STEPS} & int & Not stated & D / D \\
\code{wake\_\allowbreak termination\_\allowbreak revolutions} & \code{unsteady\_\allowbreak N\_\allowbreak revolutions\_\allowbreak wake} & \code{SET\_\allowbreak WAKE\_\allowbreak TERMINATION\_\allowbreak TIME\_\allowbreak STEPS} & float & Not stated & D / D \\
\bottomrule
\end{tabularx}}
\par\vspace{3mm}
\small \textbf{V}: verified on that build. \textbf{D}: documented, not verified on that build.
\par\vspace{2mm}
\small The command takes an integer time-step count. The revolutions key is converted using the rotor clock; it is not a second solver command. No numeric bound is declared on either model field.
\par\vspace{2mm}{\scriptsize Manual references in the database: [3, p.344].}

\setupsources{\code{cases/\_\_init\_\_.py}, \code{SolverSettings}; \code{workspace/matrix.py}, \code{\_PRESET\_ALIASES}; \code{cases/workflows.py}; \code{script/helpers.py}; \code{commands/advanced\_settings.yaml}}
\end{frame}

\begin{frame}{Wake induction and extra relaxation}
\small All commands here are inline, in phase \code{init}. Package default:
\code{None} for every key below. Floats must be finite.
\par\vspace{3mm}
{\ftstablesize\setlength{\tabcolsep}{3pt}
\begin{tabularx}{\linewidth}{@{}>{\RaggedRight\arraybackslash}p{0.18\linewidth}>{\RaggedRight\arraybackslash}p{0.18\linewidth}>{\RaggedRight\arraybackslash}p{0.26\linewidth}>{\RaggedRight\arraybackslash}p{0.07\linewidth}>{\RaggedRight\arraybackslash}p{0.10\linewidth}>{\RaggedRight\arraybackslash}X@{}}
\toprule
\textbf{Key} & \textbf{Alias} & \textbf{Solver command} & \textbf{Type} & \textbf{Solver default} & \textbf{26.123 / 26.124} \\
\midrule
\code{wake\_\allowbreak on\_\allowbreak wake\_\allowbreak induction} & None & \code{SET\_\allowbreak WAKE\_\allowbreak ON\_\allowbreak WAKE\_\allowbreak INDUCTION} & toggle & Not stated & V / D \\
\code{additional\_\allowbreak wake\_\allowbreak relaxation} & \code{additional\_\allowbreak wake\_\allowbreak relaxation\_\allowbreak iteration} & \code{ADDITIONAL\_\allowbreak WAKE\_\allowbreak RELAXATION\_\allowbreak ITERATION} & toggle & Not stated & V / D \\
\bottomrule
\end{tabularx}}
\par\vspace{3mm}
\small \textbf{V}: verified on that build. \textbf{D}: documented, not verified on that build.
\par\vspace{2mm}
\small The first switches wake-on-wake induced velocity computation. The second requests one additional wake relaxation iteration. Neither command has a stated solver default.
\par\vspace{2mm}{\scriptsize Manual references in the database: [3, pp.344-345]; [3, p.345].}

\setupsources{\code{cases/\_\_init\_\_.py}, \code{SolverSettings}; \code{workspace/matrix.py}, \code{\_PRESET\_ALIASES}; \code{cases/workflows.py}; \code{script/helpers.py}; \code{commands/advanced\_settings.yaml}}
\end{frame}

\begin{frame}{Wake termination and the row clock}
\begin{itemize}\small
  \item \code{wake\_termination\_steps} supplies the integer directly.
        For \code{wake\_termination\_revolutions}, the workflow computes
        \code{round(revolutions * steps\_per\_revolution)} and converts to an integer.
  \item Negative values count backwards from the end of the run, as the
        model and workflow describe. State one unit, not both; the rotor
        resolver refuses both together.
  \item Revolutions require a rotor speed and a clock. A steady row or
        a nonrotating unsteady row cannot convert them; use the run type's
        supported clock and the steps form where appropriate.
  \item The database describes the command as the number of time steps
        after which a wake vortex filament edge that has lost all its
        strength is removed. It states no default.
\end{itemize}
\setupsources{\code{cases/\_\_init\_\_.py}, \code{SolverSettings}; \code{workspace/matrix.py}, \code{\_PRESET\_ALIASES}; \code{cases/workflows.py}; \code{script/helpers.py}; \code{commands/advanced\_settings.yaml}}
\end{frame}

% Copyright (c) 2026 Geovana Neves. Licensed under CC BY 4.0 (https://creativecommons.org/licenses/by/4.0/). See LICENSE-AND-AUTHORSHIP.md.
%
\section{Fourteen new advanced setup keys}

\begin{frame}{Advanced setup keys (1/4)}
\relax
{\ftstablesize
\begin{tabularx}{\linewidth}{@{}>{\RaggedRight\arraybackslash}p{0.46\linewidth}>{\RaggedRight\arraybackslash}X@{}}
\toprule
\textbf{Setup key} & \textbf{Solver command} \\
\midrule
\code{laminar\_\allowbreak separation} & \code{LAMINAR\_\allowbreak SEPARATION} \\
\code{kutta\_\allowbreak joukowski\_\allowbreak lift} & \code{KUTTA\_\allowbreak JOUKOWSKI\_\allowbreak LIFT\_\allowbreak FORCES} \\
\code{aeroelastic\_\allowbreak rbf\_\allowbreak type} & \code{AEROELASTIC\_\allowbreak RBF\_\allowbreak TYPE} \\
\code{print\_\allowbreak rotor\_\allowbreak induced\_\allowbreak velocities} & \code{PRINT\_\allowbreak ROTOR\_\allowbreak INDUCED\_\allowbreak VELOCITIES} \\
\bottomrule
\end{tabularx}}
\begin{itemize}\small
\item A stated key emits its command; an omitted key emits nothing.
\item Unknown keys, unsupported values and unavailable commands are refused, naming the build.
\item The command tables below retain the per-build status. A new key does not establish licensed verification.
\end{itemize}
\end{frame}

\begin{frame}{Advanced setup keys (2/4)}
\relax
{\ftstablesize
\begin{tabularx}{\linewidth}{@{}>{\RaggedRight\arraybackslash}p{0.46\linewidth}>{\RaggedRight\arraybackslash}X@{}}
\toprule
\textbf{Setup key} & \textbf{Solver command} \\
\midrule
\code{adaptive\_\allowbreak field\_\allowbreak grid\_\allowbreak refinement} & \code{SET\_\allowbreak ADAPTIVE\_\allowbreak FIELD\_\allowbreak GRID\_\allowbreak REFINEMENT} \\
\code{rotor\_\allowbreak induced\_\allowbreak velocity\_\allowbreak blending} & \code{ROTOR\_\allowbreak INDUCED\_\allowbreak VELOCITY\_\allowbreak BLENDING} \\
\code{wake\_\allowbreak numerical\_\allowbreak relaxation} & \code{SET\_\allowbreak WAKE\_\allowbreak NUMERICAL\_\allowbreak RELAXATION} \\
\code{wake\_\allowbreak relaxation} & \code{SET\_\allowbreak WAKE\_\allowbreak RELAXATION} \\
\bottomrule
\end{tabularx}}
\begin{itemize}\small
\item A stated key emits its command; an omitted key emits nothing.
\item Unknown keys, unsupported values and unavailable commands are refused, naming the build.
\item The command tables below retain the per-build status. A new key does not establish licensed verification.
\end{itemize}
\end{frame}

\begin{frame}{Advanced setup keys (3/4)}
\relax
{\ftstablesize
\begin{tabularx}{\linewidth}{@{}>{\RaggedRight\arraybackslash}p{0.46\linewidth}>{\RaggedRight\arraybackslash}X@{}}
\toprule
\textbf{Setup key} & \textbf{Solver command} \\
\midrule
\code{wake\_\allowbreak decay\_\allowbreak constant\_\allowbreak per\_\allowbreak m} & \code{SET\_\allowbreak WAKE\_\allowbreak DECAY\_\allowbreak CONSTANT} \\
\code{wake\_\allowbreak streamwise\_\allowbreak agglomeration} & \code{SET\_\allowbreak WAKE\_\allowbreak STREAMWISE\_\allowbreak AGGLOMERATION} \\
\code{jet\_\allowbreak wake\_\allowbreak decay\_\allowbreak normalized\_\allowbreak length} & \code{SET\_\allowbreak JET\_\allowbreak WAKE\_\allowbreak DECAY\_\allowbreak NORMALIZED\_\allowbreak LENGTH} \\
\code{jet\_\allowbreak wake\_\allowbreak filaments\_\allowbreak grid\_\allowbreak induction} & \code{SET\_\allowbreak JET\_\allowbreak WAKE\_\allowbreak FILAMENTS\_\allowbreak GRID\_\allowbreak INDUCTION} \\
\bottomrule
\end{tabularx}}
\begin{itemize}\small
\item A stated key emits its command; an omitted key emits nothing.
\item Unknown keys, unsupported values and unavailable commands are refused, naming the build.
\item The command tables below retain the per-build status. A new key does not establish licensed verification.
\end{itemize}
\end{frame}

\begin{frame}{Advanced setup keys (4/4)}
\relax
{\ftstablesize
\begin{tabularx}{\linewidth}{@{}>{\RaggedRight\arraybackslash}p{0.46\linewidth}>{\RaggedRight\arraybackslash}X@{}}
\toprule
\textbf{Setup key} & \textbf{Solver command} \\
\midrule
\code{adverse\_\allowbreak gradient\_\allowbreak boundary\_\allowbreak layer} & \code{SOLVER\_\allowbreak SET\_\allowbreak ADVERSE\_\allowbreak GRADIENT\_\allowbreak BOUNDARY\_\allowbreak LAYER} \\
\code{vortex\_\allowbreak ring\_\allowbreak normalization} & \code{SOLVER\_\allowbreak VORTEX\_\allowbreak RING\_\allowbreak NORMALIZATION} \\
\bottomrule
\end{tabularx}}
\begin{itemize}\small
\item A stated key emits its command; an omitted key emits nothing.
\item Unknown keys, unsupported values and unavailable commands are refused, naming the build.
\item The command tables below retain the per-build status. A new key does not establish licensed verification.
\end{itemize}
\end{frame}


\begin{frame}{Setup toggles for separation and lift}
\small These commands have setup keys. Use the keys on a
supported build; \code{[[raw]]} remains available. Phase: \code{init}.
\par\vspace{3mm}
{\ftstablesize\setlength{\tabcolsep}{4pt}
\begin{tabularx}{\linewidth}{@{}>{\RaggedRight\arraybackslash}p{0.46\linewidth}>{\RaggedRight\arraybackslash}p{0.20\linewidth}>{\RaggedRight\arraybackslash}p{0.12\linewidth}>{\RaggedRight\arraybackslash}X@{}}
\toprule
\textbf{Command} & \textbf{Argument form} & \textbf{Solver default} & \textbf{26.123 / 26.124} \\
\midrule
\code{LAMINAR\_\allowbreak SEPARATION} & \code{ENABLE} or \code{DISABLE} & Not stated & D / D \\
\code{KUTTA\_\allowbreak JOUKOWSKI\_\allowbreak LIFT\_\allowbreak FORCES} & \code{ENABLE} or \code{DISABLE} & Not stated & D / D \\
\bottomrule
\end{tabularx}}
\par\vspace{3mm}
\small \textbf{D}: documented, not verified on that build.
\par\vspace{2mm}
\small Laminar separation is a boundary-layer separation switch. The lift command selects bound-circulation lift instead of surface-pressure integration; the database does not identify every affected output quantity.
\par\vspace{2mm}{\scriptsize Manual references: [3, p.345]; [3, p.344].}

\setupsources{\code{cases/\_\_init\_\_.py}, \code{SolverSettings}; \code{workspace/matrix.py}; \code{cases/workflows.py}, \code{\_settings}, \code{\_raw\_commands}; \code{script/solver\_setup.py}, \code{FLAG\_SPECS}; \code{commands/advanced\_settings.yaml}}
\end{frame}

\begin{frame}{Setup output and field-grid toggles}
\small These commands have setup keys. Use the keys on a
supported build; \code{[[raw]]} remains available. Phase: \code{init}.
\par\vspace{3mm}
{\ftstablesize\setlength{\tabcolsep}{4pt}
\begin{tabularx}{\linewidth}{@{}>{\RaggedRight\arraybackslash}p{0.46\linewidth}>{\RaggedRight\arraybackslash}p{0.20\linewidth}>{\RaggedRight\arraybackslash}p{0.12\linewidth}>{\RaggedRight\arraybackslash}X@{}}
\toprule
\textbf{Command} & \textbf{Argument form} & \textbf{Solver default} & \textbf{26.123 / 26.124} \\
\midrule
\code{PRINT\_\allowbreak ROTOR\_\allowbreak INDUCED\_\allowbreak VELOCITIES} & \code{ENABLE} or \code{DISABLE} & Not stated & D / D \\
\code{SET\_\allowbreak ADAPTIVE\_\allowbreak FIELD\_\allowbreak GRID\_\allowbreak REFINEMENT} & \code{ENABLE} or \code{DISABLE} & Not stated & D / D \\
\bottomrule
\end{tabularx}}
\par\vspace{3mm}
\small \textbf{D}: documented, not verified on that build.
\par\vspace{2mm}
\small The print switch writes rotor-induced velocities to the log at every unsteady time step. The refinement note marks the field-grid option as transonic only; no subsonic behavior is established.
\par\vspace{2mm}{\scriptsize Manual references: [3, p.345].}

\setupsources{\code{cases/\_\_init\_\_.py}, \code{SolverSettings}; \code{workspace/matrix.py}; \code{cases/workflows.py}, \code{\_settings}, \code{\_raw\_commands}; \code{script/solver\_setup.py}, \code{FLAG\_SPECS}; \code{commands/advanced\_settings.yaml}}
\end{frame}

\begin{frame}{Setup wake factors}
\small These commands have setup keys. Use the keys on a
supported build; \code{[[raw]]} remains available. Phase: \code{init}.
\par\vspace{3mm}
{\ftstablesize\setlength{\tabcolsep}{4pt}
\begin{tabularx}{\linewidth}{@{}>{\RaggedRight\arraybackslash}p{0.46\linewidth}>{\RaggedRight\arraybackslash}p{0.20\linewidth}>{\RaggedRight\arraybackslash}p{0.12\linewidth}>{\RaggedRight\arraybackslash}X@{}}
\toprule
\textbf{Command} & \textbf{Argument form} & \textbf{Solver default} & \textbf{26.123 / 26.124} \\
\midrule
\code{ROTOR\_\allowbreak INDUCED\_\allowbreak VELOCITY\_\allowbreak BLENDING} & \code{<float>} & 0.25 & D / D \\
\code{SET\_\allowbreak WAKE\_\allowbreak NUMERICAL\_\allowbreak RELAXATION} & \code{<float>} & 0.15 & D / D \\
\code{SET\_\allowbreak JET\_\allowbreak WAKE\_\allowbreak DECAY\_\allowbreak NORMALIZED\_\allowbreak LENGTH} & \code{<float>} & 100.0 & D / D \\
\bottomrule
\end{tabularx}}
\par\vspace{3mm}
\small \textbf{D}: documented, not verified on that build.
\par\vspace{2mm}
\small The two factors are documented from 0 to 1. Jet decay length is in jet-wake diameters, with a documented minimum of 1.0: it is the distance to one tenth of the initial strength. These notes do not enforce the bounds.
\par\vspace{2mm}{\scriptsize Manual references: [3, p.345]; [3, p.346].}

\setupsources{\code{cases/\_\_init\_\_.py}, \code{SolverSettings}; \code{workspace/matrix.py}; \code{cases/workflows.py}, \code{\_settings}, \code{\_raw\_commands}; \code{script/solver\_setup.py}, \code{FLAG\_SPECS}; \code{commands/advanced\_settings.yaml}}
\end{frame}

\begin{frame}{Setup wake decay and jet induction}
\small These commands have setup keys. Use the keys on a
supported build; \code{[[raw]]} remains available. Phase: \code{init}.
\par\vspace{3mm}
{\ftstablesize\setlength{\tabcolsep}{4pt}
\begin{tabularx}{\linewidth}{@{}>{\RaggedRight\arraybackslash}p{0.46\linewidth}>{\RaggedRight\arraybackslash}p{0.20\linewidth}>{\RaggedRight\arraybackslash}p{0.12\linewidth}>{\RaggedRight\arraybackslash}X@{}}
\toprule
\textbf{Command} & \textbf{Argument form} & \textbf{Solver default} & \textbf{26.123 / 26.124} \\
\midrule
\code{SET\_\allowbreak WAKE\_\allowbreak DECAY\_\allowbreak CONSTANT} & \code{<float>} & Not stated & D / D \\
\code{SET\_\allowbreak JET\_\allowbreak WAKE\_\allowbreak FILAMENTS\_\allowbreak GRID\_\allowbreak INDUCTION} & \code{ENABLE} or \code{DISABLE} & Not stated & Refused / refused \\
\bottomrule
\end{tabularx}}
\par\vspace{3mm}
\small \textbf{D}: documented, not verified on that build.
\par\vspace{2mm}
\small The decay constant is in inverse metres, derived in the database from \code{19.1/L}, with \code{L} in metres. No fixed default is recorded. Jet-filament grid induction has evidence only on 26.101 and 26.120; it is refused on 26.123 and 26.124.
\par\vspace{2mm}{\scriptsize Manual references: [4, p.346]; [3, p.346].}

\setupsources{\code{cases/\_\_init\_\_.py}, \code{SolverSettings}; \code{workspace/matrix.py}; \code{cases/workflows.py}, \code{\_settings}, \code{\_raw\_commands}; \code{script/solver\_setup.py}, \code{FLAG\_SPECS}; \code{commands/advanced\_settings.yaml}}
\end{frame}

\begin{frame}{Build-limited wake commands}
\small These commands have setup keys. Use the keys on a
supported build; \code{[[raw]]} remains available. Phase: \code{init}.
\par\vspace{3mm}
{\ftstablesize\setlength{\tabcolsep}{4pt}
\begin{tabularx}{\linewidth}{@{}>{\RaggedRight\arraybackslash}p{0.46\linewidth}>{\RaggedRight\arraybackslash}p{0.20\linewidth}>{\RaggedRight\arraybackslash}p{0.12\linewidth}>{\RaggedRight\arraybackslash}X@{}}
\toprule
\textbf{Command} & \textbf{Argument form} & \textbf{Solver default} & \textbf{26.123 / 26.124} \\
\midrule
\code{SET\_\allowbreak WAKE\_\allowbreak RELAXATION} & \code{ENABLE} or \code{DISABLE} & Not stated & Refused / refused \\
\code{SET\_\allowbreak WAKE\_\allowbreak STREAMWISE\_\allowbreak AGGLOMERATION} & \code{ENABLE} or \code{DISABLE} & Not stated & Refused / refused \\
\bottomrule
\end{tabularx}}
\par\vspace{3mm}
\small \textbf{D}: documented, not verified on that build.
\par\vspace{2mm}
\small Both are documented on 25.000, 25.100 and 26.000. The database records their removal from later documentation. Neither is available to emit on 26.123 or 26.124; \code{[[raw]]} does not bypass the build check.
\par\vspace{2mm}{\scriptsize Manual references: [5, p.319]; [5, p.318].}

\setupsources{\code{cases/\_\_init\_\_.py}, \code{SolverSettings}; \code{workspace/matrix.py}; \code{cases/workflows.py}, \code{\_settings}, \code{\_raw\_commands}; \code{script/solver\_setup.py}, \code{FLAG\_SPECS}; \code{commands/advanced\_settings.yaml}}
\end{frame}

\begin{frame}{Build-limited boundary-layer and vortex commands}
\small These commands have setup keys. Use the keys on a
supported build; \code{[[raw]]} remains available. Phase: \code{init}.
\par\vspace{3mm}
{\ftstablesize\setlength{\tabcolsep}{4pt}
\begin{tabularx}{\linewidth}{@{}>{\RaggedRight\arraybackslash}p{0.46\linewidth}>{\RaggedRight\arraybackslash}p{0.20\linewidth}>{\RaggedRight\arraybackslash}p{0.12\linewidth}>{\RaggedRight\arraybackslash}X@{}}
\toprule
\textbf{Command} & \textbf{Argument form} & \textbf{Solver default} & \textbf{26.123 / 26.124} \\
\midrule
\code{SOLVER\_\allowbreak SET\_\allowbreak ADVERSE\_\allowbreak GRADIENT\_\allowbreak BOUNDARY\_\allowbreak LAYER} & \code{ENABLE} or \code{DISABLE} & Not stated & Refused / refused \\
\code{SOLVER\_\allowbreak VORTEX\_\allowbreak RING\_\allowbreak NORMALIZATION} & \code{ENABLE} or \code{DISABLE} & Not stated & Refused / refused \\
\bottomrule
\end{tabularx}}
\par\vspace{3mm}
\small \textbf{D}: documented, not verified on that build.
\par\vspace{2mm}
\small Both are documented on 25.000, 25.100 and 26.000 and refused on 26.123 and 26.124. The database gives no definition of the vortex normalization beyond the switch, so no numerical mechanism is asserted here.
\par\vspace{2mm}{\scriptsize Manual references: [5, p.318].}

\setupsources{\code{cases/\_\_init\_\_.py}, \code{SolverSettings}; \code{workspace/matrix.py}; \code{cases/workflows.py}, \code{\_settings}, \code{\_raw\_commands}; \code{script/solver\_setup.py}, \code{FLAG\_SPECS}; \code{commands/advanced\_settings.yaml}}
\end{frame}

\begin{frame}{Setup aeroelastic RBF choice}
\small \code{AEROELASTIC\_RBF\_TYPE <rbf\_type>} is an inline
\code{init} command. The setup key is \code{aeroelastic\_rbf\_type}.
Its argument is one of \code{WENDLAND\_C2}, \code{GAUSSIAN},
\code{THIN\_PLATE\_SPLINE}, \code{MULTI\_QUADRATIC} or
\code{INV\_MULTI\_QUADRATIC}.
\par\vspace{1mm}
\small Solver default: \textbf{not stated}. Status: \textbf{documented} on
both 26.123 and 26.124 [3, pp.345-346]. The manual also describes numbered
choices, but the command database's accepted grammar uses the names above.
Use a quoted name as the setup value.
\begin{block}{What the choice is}
\small A radial basis function interpolates the structure's displacements,
known at its nodes, onto every aerodynamic panel [9, 10, 11]. The
multiquadric [12] has global support; the Wendland C2 [13] is compact. A
coupled FSI row states its own choice: the multiquadric for a beam line of
structural nodes (Guide 6).
\end{block}
\setupsources{\code{cases/\_\_init\_\_.py}, \code{SolverSettings}; \code{workspace/matrix.py}; \code{cases/workflows.py}, \code{\_settings}, \code{\_raw\_commands}; \code{script/solver\_setup.py}, \code{FLAG\_SPECS}; \code{commands/advanced\_settings.yaml}}
\end{frame}

% Copyright (c) 2026 Geovana Neves. Licensed under CC BY 4.0 (https://creativecommons.org/licenses/by/4.0/). See LICENSE-AND-AUTHORSHIP.md.
%
\section{Loads keys and body rates}

\begin{frame}{Loads analysis keys on steady rows}
\relax
{\ftstablesize
\begin{tabularx}{\linewidth}{@{}>{\RaggedRight\arraybackslash}p{0.20\linewidth}>{\RaggedRight\arraybackslash}p{0.30\linewidth}>{\RaggedRight\arraybackslash}X@{}}
\toprule
\textbf{Setup key} & \textbf{Solver command} & \textbf{Value} \\
\midrule
\code{analysis\_\allowbreak families} & \code{SET\_\allowbreak SOLVER\_\allowbreak ANALYSIS\_\allowbreak BOUNDARIES} & family names; every other boundary leaves the loads. Resolved like \code{vorticity\_drag\_families} \\
\code{load\_\allowbreak units} & \code{SET\_\allowbreak LOADS\_\allowbreak AND\_\allowbreak MOMENTS\_\allowbreak UNITS} & \code{COEFFICIENTS}, \code{NEWTONS}, \code{KILO-NEWTONS}, \code{POUND-FORCE}, \code{KILOGRAM-FORCE} \\
\code{inviscid\_\allowbreak loads} & \code{SET\_\allowbreak INVISCID\_\allowbreak LOADS} & toggle: the loads without their viscous part \\
\bottomrule
\end{tabularx}}
\begin{itemize}\small
  \item Stated after \code{START\_SOLVER} on every point. The solver's own
        spellings (\code{set\_loads\_and\_moments\_units}, ...) are aliases.
  \item A row of an unsteady run type stating one is \textbf{refused at plan}.
  \item A point whose loads table is not in \code{COEFFICIENTS} writes
        \textbf{no product row}; its export is still collected and hashed, and
        the post names the unit.
\end{itemize}
\setupsources{\code{cases/\_\_init\_\_.py}, \code{SolverSettings};
  \code{workspace/matrix.py}, \code{\_PRESET\_ALIASES}; \code{cases/workflows.py};
  \code{commands/solver\_analysis.yaml}}
\end{frame}

\begin{frame}{Two keys 26.124 does not answer}
\relax
{\ftstablesize
\begin{tabularx}{\linewidth}{@{}>{\RaggedRight\arraybackslash}p{0.30\linewidth}>{\RaggedRight\arraybackslash}p{0.34\linewidth}>{\RaggedRight\arraybackslash}X@{}}
\toprule
\textbf{Setup key} & \textbf{Solver command} & \textbf{On 26.124} \\
\midrule
\code{vorticity\_\allowbreak lift\_\allowbreak model} & \code{SET\_\allowbreak VORTICITY\_\allowbreak LIFT\_\allowbreak MODEL} & removed: refused at plan \\
\code{unsteady\_\allowbreak viscous\_\allowbreak coupling\_\allowbreak iteration} & \code{SET\_\allowbreak UNSTEADY\_\allowbreak VISCOUS\_\allowbreak COUPLING\_\allowbreak ITERATION} & removed: refused at plan \\
\bottomrule
\end{tabularx}}
\begin{itemize}\small
  \item Both are stated before \code{INITIALIZE\_SOLVER}. The lift model is a
        toggle on every run type; the coupling step is an integer $\geq 1$,
        unsteady run types only, and emits on 25.100 and 26.000 alone.
  \item 26.124 answers both names as \emph{Unrecognized command}, word for
        word as an unknown name [2]. A \code{[[raw]]}
        \code{SET\_VORTICITY\_LIFT\_MODEL} on a 26.124 row is refused too.
  \item The lift model beside \code{kutta\_joukowski\_lift = true} plans with a
        warning.
\end{itemize}
\setupsources{\code{cases/\_\_init\_\_.py}, \code{SolverSettings};
  \code{commands/advanced\_settings.yaml}}
\end{frame}

\begin{frame}[fragile]{Body rates: a cell of the row, not the setup}
\small \code{roll\_rate}, \code{pitch\_rate} and \code{yaw\_rate} are keys of
\code{FLIGHT\_CONDITION}, in deg/s, one non-zero rate per row; the reference's
\code{[body\_axes]} says which model axis is which.
\begin{lstlisting}[style=ftsbase]
MACH:0.2, REmi:5.5, ALPHA:0, roll_rate:sweep | -40,0,40
\end{lstlisting}
\relax
{\ftstablesize
\begin{tabularx}{\linewidth}{@{}>{\RaggedRight\arraybackslash}p{0.13\linewidth}>{\RaggedRight\arraybackslash}p{0.15\linewidth}>{\RaggedRight\arraybackslash}p{0.27\linewidth}>{\RaggedRight\arraybackslash}X@{}}
\toprule
\textbf{Rate} & \textbf{Positive is} & \textbf{Emitted for +40 deg/s} & \textbf{Measured on 26.124} \\
\midrule
\code{roll\_rate}  & right wing down & \code{ROTATION <frame> X -6.667} & rolling moment opposes the roll [2] \\
\code{pitch\_rate} & nose up         & \code{ROTATION <frame> Y 6.667}  & pitching moment nose down [2] \\
\code{yaw\_rate}   & nose right      & \code{ROTATION <frame> Z -6.667} & advancing left wing gains lift [2] \\
\bottomrule
\end{tabularx}}
\begin{alertblock}{Older rows: reversed sign}
\small Roll and yaw were emitted with the opposite sign and solved at $-p$ or
$-r$. Redo such points with \code{--force-rerun}; \code{--resume} skips them.
\end{alertblock}
\setupsources{\code{cases/workflows.py}, \code{FREESTREAM\_ROTATION\_SIGN};
  \code{workspace/matrix.py}}
\end{frame}

% Copyright (c) 2026 Geovana Neves. Licensed under CC BY 4.0 (https://creativecommons.org/licenses/by/4.0/). See LICENSE-AND-AUTHORSHIP.md.
%
\section{Writing raw entries}

\begin{frame}{The raw table and its checks}
\begin{itemize}\small
  \item Each \code{[[raw]]} table contains \code{command}, the complete
        one-line command with arguments, and \code{before}, its insertion phase.
  \item Phases: \code{geometry}, \code{setup}, \code{init}, \code{exec},
        \code{analysis}, \code{export}, \code{control}.
        The line goes before the first command of the named phase;
        \code{control} places it at the script's head.
  \item At plan, the emitter checks that the command is available on the
        row's build, that its grammar and argument count and types match,
        and that the script stays in phase order. A command from a later
        phase cannot be placed before an earlier one.
  \item One-line commands only. Keyword blocks, payloads and parameter
        blocks are refused by this route. A helper keyword alone does not
        make a top-level setup key valid.
  \item Lines retain their written order. The run record's
        \code{raw\_commands} keeps \code{command}, \code{before},
        \code{setup} and \code{source}; the provenance document carries them too.
\end{itemize}
\setupsources{\code{cases/\_\_init\_\_.py}, \code{RawCommand};
  \code{cases/workflows.py}, \code{\_raw\_commands}; \code{workspace/inputs.py};
  also repository \code{docs/workspace-and-workflows.md}, raw commands section}
\end{frame}

\begin{frame}[fragile]{Setup settings and the pproc surface window}
\small In the setup, a dedicated key replaces the raw separation command:
\begin{lstlisting}[basicstyle=\ttfamily\small,numbers=none]
laminar_separation = true
\end{lstlisting}
\small In the separate pproc artifact, request the surface average:
\begin{lstlisting}[basicstyle=\ttfamily\small,numbers=none]
[time_averaging]
last_iters = 36
\end{lstlisting}
\begin{itemize}\small
  \item With 144 time steps, the package resolves the window to 109 to 144.
        Exactly one of \code{last\_iters} or \code{last\_revs} is required.
  \item Surface averaging requires build 26.122 or later. Verified on 26.122
        [2]: the bounds are time steps, inclusive, 1-based.
  \item Existing \code{[[raw]]} entries still work. Use one route per
        setting; the pproc window also records the surface's provenance.
\end{itemize}
\end{frame}

\begin{frame}{What belongs beside the settings}
\begin{itemize}\small
  \item A raw line is fixed for every row citing this setup. Row-specific
        \code{RAW} entries extend the preset: at a given phase, preset
        lines come first and row lines follow.
  \item \code{[[flags]]} declares a command whose value the row supplies.
        It reaches only \code{control}, \code{geometry} and \code{setup}.
        The advanced \code{init} commands in this deck therefore need
        \code{[[raw]]} when they have no top-level setup key.
  \item \code{[flight\_condition]} can supply fluid pins; it is consumed
        separately from \code{SolverSettings}. It does not create new solver keys.
  \item Geometry \code{[aliases]} and \code{[[frames]]} moved to the
        reference artifact. Post-processing tables belong to the
        \code{PPROC} artifact. Their old setup location is refused.
  \item Do not state a field and its alias together. The resolver refuses
        two spellings of the same setting, even when the values agree.
\end{itemize}
\setupsources{\code{workspace/inputs.py}, \code{SetupArtifact};
  \code{workspace/matrix.py}, \code{\_solver\_from\_setup};
  \code{cases/\_\_init\_\_.py}, \code{CustomFlag}; \code{cases/workflows.py}}
\end{frame}

% Copyright (c) 2026 Geovana Neves. Licensed under CC BY 4.0 (https://creativecommons.org/licenses/by/4.0/). See LICENSE-AND-AUTHORSHIP.md.
%
\section{Example and checklist}

\begin{frame}[fragile]{Complete setup example: settings}
\small \code{inputs/setups/s901.toml}, selected by \code{SET = s901}.
An illustrative unsteady setup for 26.123 or 26.124, continued next frame.
\begin{lstlisting}[basicstyle=\ttfamily\scriptsize,numbers=none]
iterations = 500
convergence = 1e-5
forced_iterations = false
convergence_iterations = 5
max_threads = 8
timeout_s = 3600.0
walltime_margin_s = 1200.0
boundary_layer = "TURBULENT"
viscous_coupling = true
solver_model = "INCOMPRESSIBLE"
wall_collision_avoidance = false
solver_stabilization = 0.25
minimum_cp = -100.0
reynolds_averaged_drag = false
reference_velocity_m_per_s = 40.0
significant_digits = 7
\end{lstlisting}
\setupsources{\code{cases/\_\_init\_\_.py}, \code{SolverSettings};
  \code{workspace/matrix.py}; \code{cases/workflows.py}; \code{script/helpers.py}}
\end{frame}

\begin{frame}[fragile]{Complete setup example: wake settings}
\small Continue the same file. Use
geometry carrying the \code{Wing} family; these are example choices.
\begin{lstlisting}[basicstyle=\ttfamily\scriptsize,numbers=none]
load_solver_initialization = false
symmetry_loads = false
vorticity_drag_families = ["Wing"]
mesh_induced_wake_velocity = true
farfield_layers = 5          # the recommended value; above 5 is refused
unsteady_pressure_and_kutta = true
wake_termination_steps = 72
wake_on_wake_induction = true
additional_wake_relaxation = false
laminar_separation = true
\end{lstlisting}
\setupsources{\code{cases/\_\_init\_\_.py}, \code{SolverSettings};
  \code{workspace/matrix.py}; \code{cases/workflows.py};
  \code{commands/advanced\_settings.yaml}; \code{commands/unsteady\_solver.yaml}}
\end{frame}

\begin{frame}{Before planning the row}
\begin{enumerate}\small
  \item Check the \code{SET} id and the file under \code{inputs/setups/}.
        Keep settings above \code{[[raw]]} tables so TOML does not put them
        inside the last raw entry.
  \item Check each key, alias, type and bound. Toggle strings are quoted;
        TOML booleans are not. Omit optional keys for \code{None}.
  \item Check the row's build and workflow. Set the physical clock on the
        row; choose one wake-termination unit. Family names must resolve
        on that geometry.
  \item Check row overrides for \code{NCPUS} and \code{SYMMETRY\_LOADS}.
        Read recorded-only warnings; they mean no solver command was emitted.
  \item Run \code{pyfs-matrix plan}. Inspect the emitted commands, including
        the minimum-Cp default and each raw line, before running the solver.
        \code{pyfs-matrix inspect-setups} shows each resolved
        value and where it came from, and \code{plan --setup-standards
        --setup-guidelines} writes the \code{s9XX} presets and their guide.
  \item After the run, inspect \code{solver\_setup} and
        \code{raw\_commands} in the run record. A documented command is
        not a measured result on that build.
\end{enumerate}
\setupsources{\code{workspace/inputs.py}; \code{workspace/matrix.py};
  \code{cases/workflows.py}; \code{script/solver\_setup.py};
  also repository \code{docs/workspace-and-workflows.md}}
\end{frame}


\begin{frame}{The physics behind the keys, and where to read it}
\relax
{\ftstablesize\renewcommand{\arraystretch}{1.12}
\begin{tabularx}{\linewidth}{@{}>{\RaggedRight\arraybackslash}p{0.30\linewidth}>{\RaggedRight\arraybackslash}X@{}}
\toprule
\textbf{Keys} & \textbf{The physics, and its origin} \\
\midrule
The solve itself & a panel method: source and doublet singularities on the
  surface, one boundary condition per panel [6, 7] \\
\code{solver\_model} & incompressible potential flow [7]; the linearised
  subsonic correction of Prandtl and Glauert [8]; the tangent-cone and modified
  Newtonian pressure laws of supersonic and hypersonic flow [14] \\
\code{unsteady\_pressure\_and\_kutta}, trailing edges & the Kutta condition at a
  sharp trailing edge, and its unsteady form [7, 19] \\
Wake relaxation, termination, decay & a force-free wake moved with the local
  velocity; a time-stepping wake that grows one row per step [7] \\
\code{boundary\_layer}, \code{viscous\_coupling} & the boundary layer and its
  transition [15], coupled to the inviscid solution [16] \\
Separation models and criteria & turbulent separation by a pressure-gradient
  criterion [17]; a maximum-lift criterion [18] \\
\code{aeroelastic\_rbf\_type} & radial basis function interpolation of a
  structure's displacements [9, 10, 11, 12, 13] \\
\code{SYMMETRY}, periodic copies & the method of images for a mirror plane,
  and copies turned about an axis [7] \\
\bottomrule
\end{tabularx}}
\end{frame}

\begin{frame}{References}
\begin{reflist}
  \bibitem{g04r01} \docref{workspace-and-workflows}{The workspace and the workflow}
  \bibitem{g04r02} \recordsref
  \bibitem{g04r03} Altair Engineering, \emph{FlightStream User Manual}, version 26.12, 2026.
  \bibitem{g04r04} Altair Engineering, \emph{FlightStream User Manual}, version 26.1 (build 26121), 2026.
  \bibitem{g04r05} Altair Engineering, \emph{FlightStream User Guide}, version 26.0, 2025.
  \bibitem{g04r06} J. L. Hess, A. M. O. Smith, ``Calculation of potential flow about
        arbitrary bodies'', \emph{Progress in Aerospace Sciences}, 8,
        pp. 1-138, 1967.
  \bibitem{g04r07} J. Katz, A. Plotkin, \emph{Low-Speed Aerodynamics}, 2nd ed.,
        Cambridge University Press, Cambridge, 2001.
  \bibitem{g04r08} H. Glauert, ``The effect of compressibility on the lift of an
        aerofoil'', \emph{Proceedings of the Royal Society of London, Series
        A}, 118(779), pp. 113-119, 1928.
  \bibitem{g04r09} A. Beckert, H. Wendland, ``Multivariate interpolation for
        fluid-structure-interaction problems using radial basis functions'',
        \emph{Aerospace Science and Technology}, 5(2), pp. 125-134, 2001.
  \bibitem{g04r10} A. de Boer, M. S. van der Schoot, H. Bijl, ``Mesh deformation based
        on radial basis function interpolation'', \emph{Computers and
        Structures}, 85(11-14), pp. 784-795, 2007.
\end{reflist}
\end{frame}

\begin{frame}{References (continued)}
\begin{reflist}
  \setcounter{enumi}{10}
  \bibitem{g04r11} T. C. S. Rendall, C. B. Allen, ``Unified fluid-structure
        interpolation and mesh motion using radial basis functions'',
        \emph{International Journal for Numerical Methods in Engineering},
        74(10), pp. 1519-1559, 2008.
  \bibitem{g04r12} R. L. Hardy, ``Multiquadric equations of topography and other
        irregular surfaces'', \emph{Journal of Geophysical Research}, 76(8),
        pp. 1905-1915, 1971.
  \bibitem{g04r13} H. Wendland, ``Piecewise polynomial, positive definite and compactly
        supported radial functions of minimal degree'', \emph{Advances in
        Computational Mathematics}, 4(1), pp. 389-396, 1995.
  \bibitem{g04r14} J. D. Anderson, \emph{Modern Compressible Flow: With Historical
        Perspective}, 3rd ed., McGraw-Hill, New York, 2003.
  \bibitem{g04r15} H. Schlichting, K. Gersten, \emph{Boundary-Layer Theory}, 9th ed.,
        Springer, Berlin, 2017.
  \bibitem{g04r16} M. Drela, M. B. Giles, ``Viscous-inviscid analysis of transonic and
        low Reynolds number airfoils'', \emph{AIAA Journal}, 25(10),
        pp. 1347-1355, 1987.
  \bibitem{g04r17} B. S. Stratford, ``The prediction of separation of the turbulent
        boundary layer'', \emph{Journal of Fluid Mechanics}, 5(1), pp. 1-16,
        1959.
  \bibitem{g04r18} W. O. Valarezo, V. D. Chin, ``Method for the prediction of wing
        maximum lift'', \emph{Journal of Aircraft}, 31(1), pp. 103-109, 1994.
  \bibitem{g04r19} W. M. Kutta, ``Auftriebskr\"afte in str\"omenden Fl\"ussigkeiten'',
        \emph{Illustrierte Aeronautische Mitteilungen}, 6, pp. 133-135, 1902.
\end{reflist}
\end{frame}

\end{document}
